\uuid{guLj}
\titre{Dérivée seconde ou plus}
\auteur{bodin ; hanani ; mzari}
\organisation{exo7}
\licence{CC BY-NC-SA 4.0}
\datecreate{2020-01-20}
\niveau{L1}
\module{Analyse}
\chapitre{Dérivabilité des fonctions réelles}
\sousChapitre{Autre}
\difficulte{2}
\theme{Dérivées, dérivée seconde ou plus}
\exo7qcmid{lille-1.derivees.011}

\contenu{
\begin{question}
\qclassification{124.00}
\qtags{motcle=dérivée seconde ou plus}
Soit $\displaystyle f(x)=x^2\mathrm{e}^x$. Quelles sont les bonnes réponses ?
\begin{answers}
    \good{$\displaystyle f''(x)=(x^2+4x+2)\mathrm{e}^x$}
    \bad{$\displaystyle f''(x)=2\mathrm{e}^x$}
    \good{Pour $n\in \Nn^*$, $\displaystyle f^{(n)}(x)=(x^2+2nx+n^2-n)\mathrm{e}^x$.}
    \bad{Pour $n\in \Nn^*$, $\displaystyle f^{(n)}(x)=(x^2+2nx+n)\mathrm{e}^x$.}
\end{answers}
\begin{explanations}
On applique la formule de Leibniz
$$\displaystyle f^{(n)}(x)=\sum _{k=0}^n\mathrm{C}_n^k(x^2)^{(k)}(\mathrm{e}^x)^{(n-k)}=[x^2+2nx+n(n-1)]\mathrm{e}^x.$$
\end{explanations}
\end{question}
}
